\documentclass[11pt]{article}

\usepackage[letterpaper,margin=1in]{geometry}
\usepackage[T1]{fontenc}
\usepackage[utf8]{inputenc}
\usepackage{lmodern}
\usepackage{microtype}
\usepackage{amsmath,amssymb,mathtools,bm}
\usepackage{enumitem}
\usepackage{xcolor}
\usepackage{booktabs}
\usepackage{array}
\usepackage{hyperref}
\usepackage[most]{tcolorbox}
\usepackage{fancyhdr}
\usepackage{titlesec}
\usepackage{parskip}

\definecolor{qftblue}{HTML}{245A8D}
\definecolor{qftgreen}{HTML}{2F6B4F}
\definecolor{qftred}{HTML}{A23B3B}
\definecolor{qftgray}{HTML}{F4F6F8}
\definecolor{qftgold}{HTML}{C08B2C}

\hypersetup{
  colorlinks=true,
  linkcolor=qftblue,
  urlcolor=qftblue,
  citecolor=qftgreen,
  pdftitle={PHYS 662 Lecture 2: From Quantum Particles to Quantum Fields},
  pdfauthor={Nima Lashkari}
}

\pagestyle{fancy}
\fancyhf{}
\lhead{PHYS 662: Quantum Field Theory I}
\rhead{Lecture 2}
\cfoot{\thepage}
\setlength{\headheight}{14pt}

\titleformat{\section}{\Large\bfseries\color{qftblue}}{\thesection.}{0.6em}{}
\titleformat{\subsection}{\large\bfseries\color{qftgreen}}{\thesubsection}{0.6em}{}
\titleformat{\subsubsection}{\normalsize\bfseries\color{qftred}}{\thesubsubsection}{0.6em}{}

\setlist[itemize]{leftmargin=2em,itemsep=0.25em,topsep=0.4em}
\setlist[enumerate]{leftmargin=2em,itemsep=0.25em,topsep=0.4em}

\newcommand{\dd}{\mathrm{d}}
\newcommand{\e}{\mathrm{e}}
\newcommand{\ii}{\mathrm{i}}
\newcommand{\cD}{\mathcal{D}}
\newcommand{\cH}{\mathcal{H}}
\newcommand{\cO}{\mathcal{O}}
\newcommand{\R}{\mathbb{R}}
\newcommand{\avg}[1]{\left\langle #1\right\rangle}
\newcommand{\comm}[2]{\left[#1,#2\right]}

\newtcolorbox{keyidea}[1][]{
  enhanced,
  breakable,
  colback=qftgray,
  colframe=qftblue,
  boxrule=0.8pt,
  arc=1.5mm,
  title={#1},
  fonttitle=\bfseries,
  coltitle=qftblue
}

\newtcolorbox{remarkbox}[1][]{
  enhanced,
  breakable,
  colback=white,
  colframe=qftgreen,
  boxrule=0.7pt,
  arc=1.5mm,
  title={#1},
  fonttitle=\bfseries,
  coltitle=qftgreen
}

\newtcolorbox{warningbox}[1][]{
  enhanced,
  breakable,
  colback=white,
  colframe=qftred,
  boxrule=0.7pt,
  arc=1.5mm,
  title={#1},
  fonttitle=\bfseries,
  coltitle=qftred
}

\begin{document}
\pagenumbering{gobble}

\begin{titlepage}
\centering
\vspace*{1.2in}
{\scshape\Large PHYS 662 --- Quantum Field Theory I\par}
\vspace{0.35in}
{\Huge\bfseries\color{qftblue} Lecture 2\par}
\vspace{0.15in}
{\LARGE\bfseries From Quantum Particles to Quantum Fields\par}
\vspace{0.12in}
{\Large\color{qftgreen} Relativity, Locality, and the Limits of the Particle Picture\par}
\vspace{0.6in}
{\large Nima Lashkari\par}
{\large Fall 2026\par}
\vfill
\begin{minipage}{0.86\textwidth}
\small
Quantum field theory is the framework in which quantum fluctuations, special
relativity, and locality coexist.  Relativistic localization makes a
fixed-particle-number description unstable, while local fields provide
observables whose spacelike commutators vanish.  Particles remain important,
but they emerge as special excitations of the field theory rather than serving
as its universally valid starting point.

\medskip
Course page: \url{https://nimalashkari.com/quantum-field-theory-2026-part-i/}
\end{minipage}
\vspace{0.5in}
\end{titlepage}

\pagenumbering{roman}
\tableofcontents
\newpage
\pagenumbering{arabic}

\section*{Roadmap}
\addcontentsline{toc}{section}{Roadmap}

Lecture 1 showed how a classical statistical field can emerge from a
many-body system after coarse-graining.  This lecture first completes that
argument, then turns to quantum fluctuations and asks why relativistic
quantum physics cannot be organized consistently around a fixed number of
particles.

The central questions are:
\begin{enumerate}
  \item How does a block variable become a continuum field?
  \item In what sense is ordinary quantum mechanics already a field theory in
  zero spatial dimensions?
  \item Why does relativistic localization force us to allow particle
  creation and annihilation?
  \item How does a local quantum field theory restore causal behavior?
  \item When does the theory admit a particle or quasiparticle
  interpretation?
\end{enumerate}

\section{Review: from spin blocks to a classical field}

The three elementary principles emphasized in Lecture 1 were
\begin{center}
\begin{tabular}{>{\bfseries}p{0.20\textwidth} p{0.67\textwidth}}
Locality & Short-range interactions and a local derivative expansion. \\
Symmetry & Constraints on the effective free energy and its observables. \\
Fluctuations & Probability distributions, correlation functions, and
coarse-graining. \\
\end{tabular}
\end{center}
Stability, expressed through a free energy bounded below, will also become
important.  In the quantum relativistic theory it becomes part of the
spectrum condition: the energy must be bounded below and the vacuum must be
stable.

\subsection{Block magnetization}

Consider the nearest-neighbor Ising model
\begin{equation}
  H[\{s_x\}]=-J\sum_{\langle xy\rangle}s_xs_y,
  \qquad s_x=\pm1,
  \qquad J>0.
  \label{eq:IsingH}
\end{equation}
Divide the lattice into blocks $B_X$, each containing $N_B$ spins, and define
the block magnetization
\begin{equation}
  m_X=\frac{1}{N_B}\sum_{x\in B_X}s_x.
  \label{eq:blockm}
\end{equation}
Although every microscopic spin is discrete, $m_X$ takes $N_B+1$ values in
$[-1,1]$.  When $N_B\gg1$, these values are closely spaced and the block
magnetization can be treated as an approximately continuous variable.

A related soft-spin description replaces $s_x=\pm1$ by a real variable
$\phi_x$ with local weight
\begin{equation}
  p(\phi_x)=\frac{1}{z_g}\exp\left[-g(\phi_x^2-1)^2\right],
  \qquad g\gg1.
  \label{eq:softspin}
\end{equation}
Large $g$ concentrates the distribution near $\phi_x=\pm1$, while retaining
a continuous integration variable.

\subsection{The gradient term and locality}

The nearest-neighbor interaction can be rewritten, up to an additive
constant, as
\begin{equation}
  -J s_xs_{x+a\hat\mu}
  =\frac{J}{2}\left(s_{x+a\hat\mu}-s_x\right)^2-J.
  \label{eq:discretegradient}
\end{equation}
For a slowly varying field,
\begin{equation}
  \phi(x+a\hat\mu)-\phi(x)
  =a\,\partial_\mu\phi(x)+\cO(a^2).
\end{equation}
Thus a local lattice interaction produces a continuum term proportional to
$(\bm\nabla\phi)^2$.  Symmetry and stability then lead to a local
Landau--Ginzburg functional of the form
\begin{equation}
  \beta F[\phi]
  =\int\dd^dx\left[
    \frac{\kappa}{2}(\bm\nabla\phi)^2
    +\frac{r}{2}\phi^2
    +\frac{u}{4!}\phi^4+\cdots
  \right],
  \qquad u>0,
  \label{eq:LG2}
\end{equation}
where the ellipsis contains higher powers and derivatives allowed by the
$\mathbb Z_2$ symmetry $\phi\mapsto-\phi$.

\subsection{Reorganizing the partition function}

The microscopic partition function can be grouped exactly according to the
values of the block variables:
\begin{align}
  Z
  &=\sum_{\{s_x\}}\e^{-\beta H[\{s_x\}]}
  \notag\\
  &=\sum_{\{m_X\}}
    \sum_{\{s_x\}\,\text{consistent with }\{m_X\}}
    \e^{-\beta H[\{s_x\}]}
  \equiv
  \sum_{\{m_X\}}\e^{-\beta F[\{m_X\}]}.
  \label{eq:constrainedF2}
\end{align}
This equation \emph{defines} the constrained free energy.  It includes both
the energy of a microscopic configuration and the entropy associated with
the many configurations that share the same coarse-grained profile.

When the block variables are effectively continuous and vary slowly on the
cutoff scale, the sum becomes a functional integral:
\begin{equation}
  \boxed{
  Z=\int\cD m\,\e^{-\beta F[m]},
  \qquad
  \mathcal P[m]=\frac{1}{Z}\e^{-\beta F[m]}.
  }
  \label{eq:statpath}
\end{equation}
The configuration space is now a space of functions $m(\bm x)$.

\begin{keyidea}[Continuum criterion]
The continuum description is useful when the correlation length is much
larger than the microscopic lattice spacing, $\xi/a\gg1$.  A block can then
contain many microscopic degrees of freedom while remaining small compared
with the distances over which the field changes appreciably.
\end{keyidea}

\newpage
\begin{warningbox}[Classical versus quantum dimensionality]
A classical statistical system on a $d$-dimensional lattice gives a
$d$-dimensional Euclidean statistical field theory.  A $d$-dimensional
quantum system, by contrast, acquires an additional imaginary-time direction
in its Euclidean path-integral representation and is related to a
$(d+1)$-dimensional Euclidean field theory.
\end{warningbox}

\section{Quantum fluctuations and quantum correlation functions}

For a quantum system, the canonical variables become noncommuting operators,
\begin{equation}
  \comm{\hat q_i}{\hat p_j}=\ii\hbar\,\delta_{ij}.
  \label{eq:CCR}
\end{equation}
A state therefore does not assign sharp simultaneous values to all
coordinates and momenta.  For an observable $\hat A$, its variance is
\begin{equation}
  (\Delta A)^2
  =\avg{\hat A^2}-\avg{\hat A}^2.
\end{equation}
These fluctuations are present even in a pure ground state; they are not
caused by ignorance of a classical microstate or by coupling to a thermal
bath.

Quantum correlation functions are expectation values of products of
operators,
\begin{equation}
  \avg{
    \hat q_{i_1}(t_1)\cdots
    \hat q_{i_n}(t_n)
    \hat p_{j_1}(t'_1)\cdots
    \hat p_{j_m}(t'_m)
  }.
  \label{eq:quantumcorr}
\end{equation}
Because the operators need not commute, their ordering is part of the
definition.  Time-ordered, retarded, advanced, and unordered correlators
answer different physical questions.

At finite temperature,
\begin{equation}
  \avg{\hat A}_\beta
  =\frac{1}{Z}\operatorname{Tr}
   \left(\e^{-\beta\hat H}\hat A\right),
  \qquad
  Z=\operatorname{Tr}\e^{-\beta\hat H},
\end{equation}
so thermal and quantum fluctuations are both present.

\begin{remarkbox}[Two kinds of fluctuations]
Thermal fluctuations arise from a probability distribution over states.
Quantum fluctuations arise from noncommuting observables in a quantum state.
They can coexist, but they are conceptually distinct.
\end{remarkbox}

\section{Ordinary quantum mechanics as a field theory in \texorpdfstring{$0+1$}{0+1} dimensions}

Consider a particle moving in $\R^d$.  Its Heisenberg-picture coordinates
\begin{equation}
  \hat q^a(t),
  \qquad a=1,\ldots,d,
\end{equation}
are $d$ operator-valued functions of time.  In this limited but useful sense,
ordinary quantum mechanics is a theory of $d$ fields in $0+1$ dimensions:
there is one time coordinate and no spatial coordinate on the base space.
The index $a$ labels the target-space coordinate; it is not a spatial
coordinate on which the field depends.

A genuine spatial quantum field has an independent operator at every point,
\begin{equation}
  \hat\phi(\bm x,t),
  \qquad
  \hat\pi(\bm x,t),
\end{equation}
with equal-time canonical commutator
\begin{equation}
  \comm{\hat\phi(\bm x,t)}{\hat\pi(\bm y,t)}
  =\ii\hbar\,\delta^{(d)}(\bm x-\bm y).
  \label{eq:fieldCCR}
\end{equation}
One way to approach this continuum is to place coupled quantum oscillators
on a lattice and send the lattice spacing to zero.  A quantum field may
therefore be viewed as the continuum limit of infinitely many coupled quantum
degrees of freedom.

\section{Why relativity destabilizes fixed particle number}

\subsection{Localization and the Compton scale}

To localize a particle within a region of size $\Delta x$, the uncertainty
principle requires momentum components of order
\begin{equation}
  \Delta p\gtrsim\frac{\hbar}{2\Delta x}.
\end{equation}
For relativistic momenta the associated energy grows as
$\Delta E\sim c\,\Delta p$.  When $\Delta x$ approaches the Compton
wavelength
\begin{equation}
  \lambda_C=\frac{\hbar}{mc},
\end{equation}
the localization energy is comparable to the rest energy.  At sufficiently
high energy, creation of additional particle--antiparticle pairs becomes
kinematically possible.  Numerical factors depend on the localization
protocol; the invariant lesson is that a one-particle description cannot
remain closed under arbitrarily sharp localization.

\begin{keyidea}[First obstruction]
Relativity converts large momentum uncertainty into enough energy to change
particle number.  A Hilbert space containing exactly one particle is
therefore not a complete arena for relativistic local experiments.
\end{keyidea}

\subsection{The positive-energy propagation kernel}

For a free relativistic particle, the positive-energy Hamiltonian is
\begin{equation}
  H(\bm p)=\sqrt{\bm p^2c^2+m^2c^4}.
\end{equation}
The formal position-space kernel is
\begin{equation}
  K_+(t,\bm r)
  =\int\frac{\dd^d p}{(2\pi\hbar)^d}
  \exp\left[
    \frac{\ii}{\hbar}\bm p\cdot\bm r
    -\frac{\ii}{\hbar}H(\bm p)t
  \right].
  \label{eq:positivekernel}
\end{equation}
Restricting to positive energies makes this kernel nonzero outside the light
cone.  In $3+1$ dimensions and at spacelike separation
$r^2-c^2t^2>0$, its large-distance behavior contains the exponential
\begin{equation}
  K_+(t,\bm r)
  \sim
  \exp\left[
    -\frac{mc}{\hbar}\sqrt{r^2-c^2t^2}
  \right]
\end{equation}
times a power-law prefactor.  The effect is exponentially small but not
identically zero.

This does not by itself describe a superluminal signal: position states and
the positive-frequency kernel are not local observables of a relativistic
theory.  It does show why a sharply localized, positive-energy,
single-particle wavefunction is not the right fundamental object.

\subsection{Local fields and microcausality}

For a free real scalar field,
\begin{equation}
  \hat\phi(x)
  =\int\frac{\dd^d p}{(2\pi)^d}
   \frac{1}{\sqrt{2E_{\bm p}}}
  \left(
    \hat a_{\bm p}\e^{-\ii p\cdot x}
    +\hat a^\dagger_{\bm p}\e^{+\ii p\cdot x}
  \right),
  \qquad
  E_{\bm p}=\sqrt{\bm p^2+m^2},
  \label{eq:modeexpansion}
\end{equation}
where this equation uses $\hbar=c=1$.  Its commutator is
\begin{equation}
  \comm{\hat\phi(x)}{\hat\phi(y)}
  =\ii\,\Delta(x-y).
  \label{eq:PauliJordan}
\end{equation}
The positive- and negative-frequency parts cancel at spacelike separation,
so
\begin{equation}
  \boxed{
  \comm{\hat\phi(x)}{\hat\phi(y)}=0
  \quad\text{when}\quad
  (x-y)^2<0.
  }
  \label{eq:microcausality}
\end{equation}
This is microcausality: operations performed in spacelike separated regions
cannot influence one another through local observables.

\begin{warningbox}[What performs the cancellation?]
It is the full local field commutator, including both frequency signs, that
vanishes outside the light cone.  It is not generally correct to say that a
particle probability is canceled by an antiparticle probability.  A real
scalar particle is its own antiparticle, yet its field is perfectly causal.
\end{warningbox}

\section{Fock space and particle-number-changing processes}

The natural weak-coupling state space includes every particle number:
\begin{equation}
  \cH_{\mathrm F}
  =
  \bigoplus_{n=0}^{\infty}\cH_n,
  \qquad
  \cH_n=
  \begin{cases}
    \operatorname{Sym}\cH_1^{\otimes n}, & \text{bosons},\\
    \operatorname{Alt}\cH_1^{\otimes n}, & \text{fermions}.
  \end{cases}
  \label{eq:Fock}
\end{equation}
Creation and annihilation operators connect neighboring sectors.  For
several particle species, the corresponding Fock spaces are combined.
Interactions may conserve a charge while changing the total number of
quanta.

Fock space is indispensable for free and perturbative QFT, but it is not the
definition of QFT.  In strongly interacting theories, curved spacetime, or
gauge theories, a preferred particle basis may be absent or only
approximate.  Local operators and their correlation functions are more
general.

\subsection{Quantum fluctuations and pair creation}

Quantum fields fluctuate even in the vacuum:
\begin{equation}
  \avg{0|\hat\phi(x)\hat\phi(y)|0}\neq0.
\end{equation}
Interactions or external backgrounds can convert available energy into
on-shell particle--antiparticle pairs.  However, a stable time-independent
vacuum does not continually emit real particles: energy and all conserved
charges must still be respected.  Vacuum correlation functions and virtual
intermediate states should not be confused with directly observable
particles.

\begin{remarkbox}[Classical and quantum parallels]
In classical statistical field theory, thermal fluctuations lead to a
probability distribution on field configurations.  In QFT, quantum
fluctuations appear in noncommuting field operators and their correlation
functions.  Both settings use generating functionals and functional
integrals, but the physical interpretation is different.
\end{remarkbox}

\section{Particles and quasiparticles are regime-dependent}

\subsection{Short-range interactions and asymptotic states}

Suppose two localized wave packets are separated by a distance $R$.  If the
interaction energy vanishes sufficiently rapidly as $R\to\infty$, the
Hamiltonian approximately splits,
\begin{equation}
  H\simeq H_1+H_2,
\end{equation}
and the state approximately factorizes into two independently propagating
excitations.  Bringing the packets together produces scattering; taking
them far apart again defines asymptotic in- and out-particles.

In an interacting theory, a stable particle is visible as an isolated pole
of an appropriate two-point function.  Schematically,
\begin{equation}
  G(p)
  \sim\frac{Z_{\mathrm{pole}}}
  {p^2-m_{\mathrm{phys}}^2+\ii\epsilon}
  +\text{continuum}.
  \label{eq:particlepole}
\end{equation}
The pole location gives the physical mass and its residue measures the
overlap of the field with the one-particle state.  An unstable excitation
produces a resonance rather than a pole on the physical sheet.

\subsection{Strong coupling, long-range forces, and confinement}

If interactions do not become negligible at large separation, the ordinary
particle picture can fail.  Confinement is the sharp example.  In pure
Yang--Mills theory, the energy of a separated color source and antisource
grows approximately as
\begin{equation}
  V(R)\simeq \sigma R
\end{equation}
at large $R$, with string tension $\sigma$.  Isolated colored asymptotic
states are then absent; the spectrum is organized into color-neutral bound
states.

In full QCD with dynamical quarks, the flux tube can break by pair creation,
so the static potential eventually saturates.  The conclusion remains that
quarks are not isolated asymptotic particles.  Hadrons, rather than quarks,
provide the particle states measured at low energy.

Long-range interactions require additional care even without confinement.
In QED, for example, a charged particle is accompanied by a soft photon
cloud, modifying the naive Fock-particle picture.

\begin{keyidea}[The particle concept is asymptotic]
Particles and quasiparticles are exceptionally useful when the theory has
stable, well-separated excitations.  They are properties of a regime and of
the spectrum, not universal microscopic ingredients.  Some strongly
coupled theories have only composite particles; others have no useful
particle description at all.
\end{keyidea}

\section{Structural consequences of relativistic QFT}

\subsection{Antiparticles}

Relativistic local QFT naturally contains both signs of frequency.  For a
complex field, the negative-frequency term is reinterpreted through
creation operators for antiparticles.  Particles and antiparticles have the
same mass and spin and opposite additive charges.

The modern statement relies on the full structure of local relativistic QFT,
including a stable vacuum and CPT symmetry.  Not every particle has a
\emph{distinct} antiparticle: a neutral real scalar, the photon, and other
self-conjugate excitations can be their own antiparticles.

\subsection{Identical particles and the spin--statistics connection}

Quanta created by the same field are intrinsically identical.  In
$3+1$-dimensional local relativistic QFT with positive energy,
\begin{equation}
  \begin{array}{ccl}
  \text{integer spin} &\Longrightarrow& \text{Bose statistics},\\
  \text{half-integer spin} &\Longrightarrow& \text{Fermi statistics}.
  \end{array}
  \label{eq:spinstat}
\end{equation}
This is the spin--statistics theorem.  It is not a statement of
nonrelativistic quantum mechanics alone; locality, Lorentz symmetry, and
the spectrum condition are essential assumptions.

In $2+1$ spacetime dimensions the topology of particle exchange is
different, and anyonic or more general braid statistics are possible.  The
dimensional assumptions should therefore always be stated.

\section{What is a quantum field?}

A quantum field is often introduced as an operator-valued function
$\hat\phi(x)$.  The mathematically and physically safer statement is that it
is an \emph{operator-valued distribution}.  The object at an exact point is
typically too singular to be an operator on the Hilbert space.  One defines
a smeared field
\begin{equation}
  \hat\phi(f)
  =\int\dd^{d+1}x\,f(x)\hat\phi(x),
  \label{eq:smearedfield}
\end{equation}
where $f$ is a smooth test function of compact support.  The support of $f$
specifies the spacetime region in which the observable is localized.

\subsection{From a classical equation to a quantum field}

For a classical real scalar field with potential $V(\phi)$, a standard
action is
\begin{equation}
  S[\phi]
  =\int\dd t\,\dd^d x\left[
    \frac12(\partial_t\phi)^2
    -\frac12(\bm\nabla\phi)^2
    -V(\phi)
  \right].
\end{equation}
Its Euler--Lagrange equation is
\begin{equation}
  \partial_t^2\phi-\bm\nabla^2\phi+V'(\phi)=0.
  \label{eq:classicalEOM}
\end{equation}
Quantization promotes the field and its conjugate momentum to
operator-valued distributions satisfying canonical commutators, or
equivalently constructs their correlation functions from a path integral.

For the free potential $V(\phi)=m^2\phi^2/2$, the equation becomes
\begin{equation}
  (\Box+m^2)\phi=0,
\end{equation}
in units with $c=\hbar=1$.  Each positive-frequency classical solution
labels a one-particle wave packet, while the full quantum field also creates
and annihilates arbitrary multiparticle states.

\begin{warningbox}[Fields are not ordinary pointwise observables]
Products such as $\hat\phi(x)^2$ are singular and must be defined by a
renormalization prescription.  Smearing controls the distributional
singularity of a single field, but composite operators require additional
work.
\end{warningbox}

\section{Three principles of relativistic QFT}

\subsection{Locality}

Local observables commute at spacelike separation,
\begin{equation}
  \comm{\cO_1(x)}{\cO_2(y)}=0,
  \qquad (x-y)^2<0,
\end{equation}
with the corresponding graded relation for fermionic fields.  This is the
operational form of relativistic causality.

\subsection{Symmetry}

Spacetime and internal symmetries act on fields, states, and observables.
Continuous symmetries lead to conserved currents and charges.  Gauge
symmetry is more subtle: it is a redundancy of description whose constraints
organize the physical state space and the local observables.

\subsection{The spectrum condition}

The Hamiltonian is bounded below, and in a Poincare-invariant theory the
joint energy--momentum spectrum lies in the closed future light cone:
\begin{equation}
  p^0\ge0,
  \qquad
  p_\mu p^\mu\ge0.
\end{equation}
This condition expresses stability and selects the vacuum as a
lowest-energy state.

\begin{keyidea}[A working definition]
A relativistic QFT assigns quantum observables to spacetime regions in a way
compatible with locality, symmetry, and a positive-energy spectrum.
Fields are the local operator-valued distributions used to generate and
organize those observables.
\end{keyidea}

\section{Two complementary quantization viewpoints}

\begin{description}
  \item[Canonical or second quantization.]
  Begin with classical fields and their conjugate momenta, impose canonical
  commutation or anticommutation relations, and construct a Hilbert space.
  For free fields this leads directly to creation and annihilation
  operators and a Fock-space particle interpretation.

  \item[Path-integral quantization.]
  Begin with the action and sum over field histories with weight
  $\e^{\ii S/\hbar}$ in Lorentzian signature, or $\e^{-S_E/\hbar}$ after
  suitable continuation to Euclidean time.  This approach keeps Lagrangian
  symmetries manifest and is well suited to perturbative and
  nonperturbative methods.
\end{description}

The two viewpoints encode the same theory when both are well defined, but
they illuminate different structures.  Neither requires us to regard
particles as fundamental.

\section{Questions to think about}

\begin{enumerate}
  \item Why is a nonzero positive-frequency kernel outside the light cone
  not, by itself, a measurable violation of causality?
  \item Which cancellation makes the free scalar commutator vanish at
  spacelike separation?
  \item What distinguishes a virtual contribution to a correlator from an
  on-shell particle detected in an experiment?
  \item What analytic feature of a two-point function identifies a stable
  particle?
  \item Can a local QFT have no particle interpretation?  What remains
  meaningful in that case?
  \item Why may a neutral particle be its own antiparticle?
  \item Which assumptions enter the spin--statistics theorem, and why do
  lower-dimensional systems admit exceptions?
\end{enumerate}

\section*{Summary}
\addcontentsline{toc}{section}{Summary}

Coarse-graining a many-body system reorganizes the microscopic partition sum
into a functional integral over smooth fields.  Quantization replaces
classical random variables by noncommuting operator-valued distributions.
Special relativity then makes a fixed-particle-number description
inadequate: localization reaches the pair-production scale, and
positive-energy one-particle propagation alone does not furnish local causal
observables.  A full local field restores microcausality through the
cancellation of its positive- and negative-frequency parts.  Fock space
describes variable particle number when a particle basis is useful, while
correlation functions and local observables remain meaningful even at strong
coupling or in regimes without particles.  Locality, symmetry, and the
positive-energy spectrum form the structural core of relativistic QFT.

\end{document}
